% GENERATED FILE. Edit canonical lessons, then run npm run generate:books.
% canonical-source-hash: fnv1a64:07ec1ca1d44a715e
% canonical-lessons: SA-C200-pathanam, SA-C200-sravanam, SA-C200-bhasanam, SA-C200-carca, SA-C200-vivadah

\chapter{Reading, Listening, Speech, Discussion, Dispute}
\label{ch:sa-reading-listening-speech-discussion-dispute}
\hlchaptermodality{200}

\begin{chapteropening}
\textbf{By the end of this chapter:} \emph{I can say reading, listening, a speech, a discussion and a dispute.}

Complete the last lesson of chapter 200: I can say reading, listening, a speech, a discussion and a dispute.
\end{chapteropening}

\section[\texorpdfstring{\sk{पठनम्} (paṭhanam)}{पठनम् (paṭhanam)}]{\sk{पठनम्} (paṭhanam) — reading}

\label{lesson:SA-C200-pathanam}



\begin{quote}
Before the new one: say the Sanskrit for mathematics, then the Sanskrit for science.
\end{quote}

\subsection*{You'll want to know: \sk{पठनम्}}
\textbf{\sk{पठनम्}} — \emph{paṭhanam} — \textquotedblleft{}reading\textquotedblright{}.

\begin{grammarlens}[title={What you've built}]
One more word for this chapter.
\end{grammarlens}

\subsection*{Your turn}
\begin{itemize}
\raggedright
  \item \emph{Say it:} \emph{paṭhanam}
  \item \emph{Say it:} \emph{paṭhanam}, once more
  \item \emph{Say it:} say \emph{paṭhanam}
  \item \emph{Recall:} say the Sanskrit for mathematics, then the Sanskrit for science, then say \emph{paṭhanam} again
\end{itemize}

\begin{tcolorbox}[breakable,colback=teal!4,colframe=teal!35!black,title={Before you move on}]
What does \sk{पठनम्} mean? (\textquotedblleft{}Reading\textquotedblright{}.) Say it once more.
\end{tcolorbox}
\section[\texorpdfstring{\sk{श्रवणम्} (śravaṇam)}{श्रवणम् (śravaṇam)}]{\sk{श्रवणम्} (śravaṇam) — listening}

\label{lesson:SA-C200-sravanam}



\begin{quote}
Before the new one: say the Sanskrit for science, then the Sanskrit for reading.
\end{quote}

\subsection*{You'll want to know: \sk{श्रवणम्}}
\textbf{\sk{श्रवणम्}} — \emph{śravaṇam} — \textquotedblleft{}listening\textquotedblright{}.

\begin{grammarlens}[title={What you've built}]
One more word for this chapter.
\end{grammarlens}

\subsection*{Your turn}
\begin{itemize}
\raggedright
  \item \emph{Say it:} \emph{śravaṇam}
  \item \emph{Say it:} \emph{śravaṇam}, once more
  \item \emph{Say it:} say \emph{śravaṇam}
  \item \emph{Recall:} say the Sanskrit for science, then the Sanskrit for reading, then say \emph{śravaṇam} again
\end{itemize}

\begin{tcolorbox}[breakable,colback=teal!4,colframe=teal!35!black,title={Before you move on}]
What does \sk{श्रवणम्} mean? (\textquotedblleft{}Listening\textquotedblright{}.) Say it once more.
\end{tcolorbox}
\section[\texorpdfstring{\sk{भाषणम्} (bhāṣaṇam)}{भाषणम् (bhāṣaṇam)}]{\sk{भाषणम्} (bhāṣaṇam) — a speech}

\label{lesson:SA-C200-bhasanam}



\begin{quote}
Before the new one: say the Sanskrit for reading, then the Sanskrit for listening.
\end{quote}

\subsection*{You'll want to know: \sk{भाषणम्}}
\textbf{\sk{भाषणम्}} — \emph{bhāṣaṇam} — \textquotedblleft{}a speech\textquotedblright{}.

\begin{grammarlens}[title={What you've built}]
One more word for this chapter.
\end{grammarlens}

\subsection*{Your turn}
\begin{itemize}
\raggedright
  \item \emph{Say it:} \emph{bhāṣaṇam}
  \item \emph{Say it:} \emph{bhāṣaṇam}, once more
  \item \emph{Say it:} say \emph{bhāṣaṇam}
  \item \emph{Recall:} say the Sanskrit for reading, then the Sanskrit for listening, then say \emph{bhāṣaṇam} again
\end{itemize}

\begin{tcolorbox}[breakable,colback=teal!4,colframe=teal!35!black,title={Before you move on}]
What does \sk{भाषणम्} mean? (\textquotedblleft{}A speech\textquotedblright{}.) Say it once more.
\end{tcolorbox}
\section[\texorpdfstring{\sk{चर्चा} (carcā)}{चर्चा (carcā)}]{\sk{चर्चा} (carcā) — a discussion}

\label{lesson:SA-C200-carca}



\begin{quote}
Before the new one: say the Sanskrit for listening, then the Sanskrit for a speech.
\end{quote}

\subsection*{You'll want to know: \sk{चर्चा}}
\textbf{\sk{चर्चा}} — \emph{carcā} — \textquotedblleft{}a discussion\textquotedblright{}.

\begin{grammarlens}[title={What you've built}]
One more word for this chapter.
\end{grammarlens}

\subsection*{Your turn}
\begin{itemize}
\raggedright
  \item \emph{Say it:} \emph{carcā}
  \item \emph{Say it:} \emph{carcā}, once more
  \item \emph{Say it:} say \emph{carcā}
  \item \emph{Recall:} say the Sanskrit for listening, then the Sanskrit for a speech, then say \emph{carcā} again
\end{itemize}

\begin{tcolorbox}[breakable,colback=teal!4,colframe=teal!35!black,title={Before you move on}]
What does \sk{चर्चा} mean? (\textquotedblleft{}A discussion\textquotedblright{}.) Say it once more.
\end{tcolorbox}
\section[\texorpdfstring{\sk{विवादः} (vivādaḥ)}{विवादः (vivādaḥ)}]{\sk{विवादः} (vivādaḥ) — a dispute}

\label{lesson:SA-C200-vivadah}



\begin{quote}
Before the new one: say the Sanskrit for a speech, then the Sanskrit for a discussion.
\end{quote}

\subsection*{You'll want to know: \sk{विवादः}}
\textbf{\sk{विवादः}} — \emph{vivādaḥ} — \textquotedblleft{}a dispute\textquotedblright{}.

\begin{grammarlens}[title={What you've built}]
One more word for this chapter.
\end{grammarlens}

\subsection*{Your turn}
\begin{itemize}
\raggedright
  \item \emph{Say it:} \emph{vivādaḥ}
  \item \emph{Say it:} \emph{vivādaḥ}, once more
  \item \emph{Say it:} say \emph{vivādaḥ}
  \item \emph{Recall:} say the Sanskrit for a speech, then the Sanskrit for a discussion, then say \emph{vivādaḥ} again
\end{itemize}

\begin{tcolorbox}[breakable,colback=teal!4,colframe=teal!35!black,title={Before you move on}]
What does \sk{विवादः} mean? (\textquotedblleft{}A dispute\textquotedblright{}.) Say it once more.
\end{tcolorbox}
